Společenské tance (často označované anglickým termínem \emph{social dance}) jsou tance, které se tančí především za účelem zábavy a společenské interakce. V~posledních letech nabývají na popularitě a s tím roste i množství tanečních kurzů, workshopů a akcí. Tanců je několik, avšak s ohledem na cílovou skupinu se tato práce zaměřuje na salsu, bachatu a zouk.

S narůstajícím počtem událostí se sice rozšiřují možnosti pro taneční komunitu, zároveň je ale stále náročnější se v této nabídce přehledně orientovat. Pro začínající tanečníky bývá obtížné zjistit, kde a kdy se konají vhodné akce a jaké jsou jejich podmínky (úroveň, typ události, cena či lokalita). Zkušenější tanečníci naopak často hledají konkrétnější akce podle stylu nebo místa a potřebují rychle porovnat více možností. Z pohledu organizátorů je klíčová především přehledná správa událostí a registrací účastníků.

V praxi jsou informace o událostech často rozptýlené napříč více kanály. Část organizátorů zveřejňuje akce na webových stránkách nebo v kalendářích tanečních škol, významná část komunikace probíhá na sociálních sítích (typicky Facebook a Instagram) a registrace bývají řešeny externě, například přes Google Formuláře či zprávy. Tato roztříštěnost vede k tomu, že tanečníci musí události složitě dohledávat, a i po registraci mohou postrádat jasné potvrzení nebo přehled o svých přihláškách. Organizátorům naopak chybí jednotné místo pro prezentaci akcí a nástroje pro přehlednou administraci účastníků.

Tuto problematiku řeší tato diplomová práce, jejímž tématem je návrh a implementace webové aplikace pro správu společenských tanečních akcí a kurzů, která by spojovala tanečníky a organizátory. Nápad na téma vznikl z osobního zájmu autorky o společenské tance a z pozorování, že existuje mezera na trhu pro platformu, která by poskytla centralizovaný a přehledný přístup k událostem i jejich správě.

Cílem této práce je navrhnout, implementovat a vhodně otestovat webovou aplikaci pro správu společenských tanečních akcí a kurzů, která tanečníkům usnadní vyhledávání událostí a organizátorům umožní přehlednou správu akcí a registrací účastníků.